\chapter{Talking to Hardware}
\label{ch:hardware}

A peripheral register is a fixed memory location with side effects; a structure
read off storage or the wire is a foreign byte layout. Ada lets the \emph{type}
carry both -- the exact address, the bit positions, the volatility, the
endianness -- so once a register or structure is described correctly it is safe at
every use, with no pointer casts or shift-and-mask macros. This chapter is those
techniques; they build directly on the language features of
Chapter~\ref{ch:embedded}.

\section{Register Addressing in Ada}
\label{ch:registers}\index{register addressing}

A peripheral register is a fixed memory location with side effects. To program it
in Ada you describe \emph{three} things: its layout (a type), its location (an
address), and the fact that touching it is observable (volatility). You never
need a pointer cast.

\subsection{An object at a fixed address}

The simplest register is a 32-bit word. Place an Ada object at its address and
mark it \texttt{Volatile} so the compiler never caches, reorders, or elides
accesses:

\begin{lstlisting}
with System;
with Interfaces; use Interfaces;

procedure Toggle_Raw is
   --  The GPIO "output set" and "output clear" W1TS/W1TC registers.
   GPIO_OUT_W1TS : Unsigned_32
     with Volatile, Address => System'To_Address (16#6000_4008#);
   GPIO_OUT_W1TC : Unsigned_32
     with Volatile, Address => System'To_Address (16#6000_400C#);
begin
   GPIO_OUT_W1TS := 2 ** 2;   --  set GPIO2  (writing a 1 sets that bit)
   GPIO_OUT_W1TC := 2 ** 2;   --  clear GPIO2 (writing a 1 clears it)
end Toggle_Raw;
\end{lstlisting}

\texttt{System'To\_Address} turns a numeric literal into an address; the
\texttt{Address} aspect binds the object there. Because the type is exactly
\texttt{Unsigned\_32}, the write is exactly one 32-bit store -- the compiler
cannot ``optimise'' it into a byte write, which a peripheral might not accept.

\cnote{this replaces \texttt{*(volatile uint32\_t *)0x60004008 = x;}\ -- but with
no cast, and with the access width fixed by the type, so the compiler cannot
silently split or merge the access.}

\subsection{Generated register packages}

Writing every register by hand is error-prone, so the whole ESP32-S3 register map
is generated from the vendor SVD file into typed packages under
\texttt{ESP32S3\_Registers.*} (one package per peripheral). Most registers are
records of named bit-fields with a representation clause; a few -- like the
write-one-to-set ports -- are bare 32-bit words. Each peripheral is one big record
placed at its base address, and a driver writes them by name:

\begin{lstlisting}
with ESP32S3_Registers.GPIO; use ESP32S3_Registers.GPIO;

--  GPIO_Periph is the peripheral record, already placed at 16#6000_4000#.
GPIO_Periph.OUT_W1TS := 2 ** 2;          --  set GPIO2 high (write-1-to-set)
--  ... and a register that IS a field record is written field by field:
GPIO_Periph.PIN (2).INT_ENA := 1;        --  GPIO2's interrupt enable, by name
\end{lstlisting}

The application almost never touches these directly -- it calls the hand-written
HAL (\texttt{ESP32S3.GPIO} and friends) -- but the same techniques underpin every
driver.

\subsection{Array overlays}

Sometimes the SVD flattens a regular bank of registers (a FIFO, a descriptor
table) into separate fields. You can overlay your own array on top of them with
\texttt{Address}, recovering clean indexing:

\begin{lstlisting}
type Reg_Bank is array (0 .. 7) of Interfaces.Unsigned_32 with Volatile;
Bank : Reg_Bank with Import, Volatile, Address => Some_Periph.First_Reg'Address;
--  Bank (3) is now the fourth 32-bit register, addressed normally.
\end{lstlisting}

The \texttt{Import} aspect says ``this object already exists at that address; do
not allocate or initialise it.'' This idiom appears throughout the HAL where the
register layout is regular but the generator did not capture it as an array.

\section{Bit Fields and Record Representation}
\label{ch:bitfields}\index{representation clause}

Hardware registers pack many fields into one word: a 3-bit mode here, a single
enable bit there, a 12-bit divider somewhere else. In C this means shift-and-mask
macros that are easy to get wrong. In Ada you declare a record whose components
have the right types and then a \emph{representation clause} that nails each
component to specific bits. The compiler does the shifting and masking for you,
correctly, every time.

\subsection{A register as a record}

\begin{lstlisting}
type Mode_Bits is mod 2 ** 3;      --  a 3-bit field
type Divider   is mod 2 ** 12;     --  a 12-bit field

type Control_Register is record
   Enable   : Boolean;             --  1 bit
   Mode     : Mode_Bits;           --  3 bits
   Reserved : Boolean := False;
   Div      : Divider;             --  12 bits
end record
  with Volatile_Full_Access, Size => 32, Bit_Order => System.Low_Order_First;

for Control_Register use record
   Enable   at 0 range 0  .. 0;
   Mode     at 0 range 1  .. 3;
   Reserved at 0 range 4  .. 4;
   Div      at 0 range 5  .. 16;
end record;
\end{lstlisting}

Now you assign fields by name and the layout is guaranteed:

\begin{lstlisting}
Ctrl := (Enable => True, Mode => 2, Div => 1023, others => <>);
\end{lstlisting}

\texttt{at 0 range L\,..\,H} means ``bits $L$ through $H$ of the byte at offset
0.'' \texttt{Volatile\_Full\_Access} guarantees the whole 32-bit word is read or
written in one access -- essential for registers that latch on a full-word
write. This is precisely how the generated \texttt{ESP32S3\_Registers} packages
describe every register, and how the HAL builds register values:

\begin{lstlisting}
--  From the real SPI driver: build a USER register value by field, in one write.
Regs.USER := (DOUTDIN => True, USR_MOSI => True, USR_MISO => True,
              CK_OUT_EDGE => Out_Edge, others => <>);
\end{lstlisting}

\subsection{Enumerations with explicit codes}

When a field's values are an enumeration, give the enumeration an explicit
representation so the symbolic names map to the exact hardware codes:

\begin{lstlisting}
type Parity is (None, Even, Odd);
for Parity use (None => 0, Even => 2, Odd => 3);
\end{lstlisting}

\subsection{Scalar storage order: decoding on-disk and on-wire data}

The same machinery decodes \emph{external} byte layouts, not just registers. A
record can declare its \texttt{Scalar\_Storage\_Order}, and GNAT will read each
field with the correct endianness automatically -- invaluable for filesystem
and protocol structures. The ext4 filesystem in Chapter~\ref{ch:storage} relies
on exactly this idea (though it ultimately uses explicit little-/big-endian
readers for clarity), and JBD2's big-endian journal records show why being able
to say ``this structure is big-endian'' in the type is so useful.

\section{Volatile, Atomic, and full-word access}
\label{sec:volatile}\index{Volatile}

Placing an object at a register's address is only half the story; you must also
tell the compiler how that location behaves, because the optimiser otherwise
assumes ordinary memory -- caching reads, dropping ``dead'' writes, reordering.
A small family of aspects states the truth:

\begin{itemize}
  \item \texttt{Volatile} -- every read and write happens, in program order, and
        none is cached or elided. The baseline for any register.
  \item \texttt{Atomic} -- additionally, each access is indivisible (a single
        load or store). Use it for a flag shared with an interrupt handler or the
        other core, where a torn half-update would be a bug.
  \item \texttt{Volatile\_Full\_Access} -- the whole object is read or written in
        \emph{one} access of its size. Essential for a register that latches on a
        full-word write or clears on read: updating one bit-field still emits a
        single 32-bit access, not a byte poke.
  \item \texttt{Volatile\_Components} / \texttt{Atomic\_Components} -- the same,
        applied to each element of an array. The runtime's per-core
        \texttt{Switch\_Pending} and \texttt{In\_Native\_Int} flags are
        \texttt{Volatile\_Components} arrays for exactly this reason.
\end{itemize}

\begin{lstlisting}
   Ready : Boolean with Atomic, Address => Status_Reg'Address;  -- ISR-shared flag
   --  Control_Register from the bit-field section is Volatile_Full_Access:
   --  Ctrl.Enable := True  emits one 32-bit store, not a read-modify-byte-write.
\end{lstlisting}

Pick the weakest that is correct: \texttt{Volatile} for most registers,
\texttt{Volatile\_Full\_Access} for latch-on-write ones, \texttt{Atomic} only
where indivisibility actually matters (it can be more expensive).

\section{Packing: \texttt{pragma Pack} and packed arrays}
\label{sec:pack}\index{packed arrays}

The bit-field record nails \emph{named} fields to bits. When you instead have a
homogeneous run of small components -- a bit vector, a table of 2-bit codes --
\texttt{pragma Pack} (or a \texttt{Component\_Size} clause) tells the compiler to
store them at their natural bit width instead of rounding each up to a byte or
word:

\begin{lstlisting}
   type Pin_Mask is array (0 .. 31) of Boolean;
   pragma Pack (Pin_Mask);          -- 32 bits total, one bit per pin
   --  or, explicitly:
   type Codes is array (0 .. 15) of Mode_Bits with Component_Size => 2;  -- 32 bits
\end{lstlisting}

A packed \texttt{array (...) of Boolean} is the idiomatic Ada bitmap: you index
it by position (\texttt{Mask (Pin) := True}) and the compiler does the
shift-and-mask, while the whole thing still occupies one word you can overlay on
a register. Reach for the named record when the fields differ; reach for a packed
array when they are uniform and indexed.

\section{Placing objects in memory}
\label{sec:placement}\index{Linker\_Section}

On a microcontroller \emph{where} an object lives is part of its meaning: some
data must sit in fast IRAM, some in external PSRAM, some in RTC memory that
survives deep sleep, and some must \emph{not} be cleared at startup. Two aspects
cover it:

\begin{itemize}
  \item \texttt{Linker\_Section => "..."} places an object in a named output
        section the linker script positions (e.g.\ an IRAM or \texttt{.noinit}
        section).
  \item an \texttt{Address} overlay (\S\ref{ch:registers}) pins it at a fixed
        location -- which is how the RTC-retained boot counter in the deep-sleep
        example lives in RTC slow memory and survives a reset:
\end{itemize}

\begin{lstlisting}
   Boot_Count : Unsigned_32 with Volatile,
     Address => System'To_Address (16#5000_0000#);   -- RTC slow RAM, retained

   IRAM_Buf : Buffer with Linker_Section => ".iram1.data";
\end{lstlisting}

Alignment matters too: \texttt{with Alignment => 64} forces DMA buffers onto a
cache-line boundary, and the GDMA driver relies on aligned, internal-RAM (not
PSRAM) descriptors -- placement and alignment together.

\begin{table}[ht]
\centering
\begin{tabular}{@{}l l l@{}}
\toprule
\textbf{Base} & \textbf{Top} & \textbf{Region} \\ \midrule
\texttt{16\#3C00\_0000\#} & \texttt{16\#3D00\_0000\#} & Flash $\to$ DROM -- cached read-only data \\
\texttt{16\#3D00\_0000\#} & \texttt{16\#3E00\_0000\#} & External PSRAM -- mapped by the bootloader \\
\texttt{16\#3FC8\_8000\#} & \texttt{16\#3FD0\_0000\#} & Internal SRAM (DRAM view) -- data, stacks, heap \\
\texttt{16\#4037\_0000\#} & \texttt{16\#403E\_0000\#} & Internal SRAM (IRAM view) -- code, vectors, DMA \\
\texttt{16\#4200\_0000\#} & \texttt{16\#4400\_0000\#} & Flash $\to$ IROM -- cached code \\
\texttt{16\#5000\_0000\#} & \texttt{16\#5000\_2000\#} & RTC slow RAM -- retained across deep sleep \\
\bottomrule
\end{tabular}
\caption{Principal address-space regions, by ascending base address. Flash
(DROM) and external PSRAM are not separate windows: they are two sub-ranges of
one 32\,MB external data-bus window (\texttt{16\#3C00\_0000\#}--\texttt{16\#3E00%
\_0000\#}) that the cache MMU partitions -- this project maps flash from the
bottom and PSRAM at \texttt{16\#3D00\_0000\#}, the PSRAM's actual extent set by
the fitted device. Placement pins objects where they belong: RTC memory for
state that survives sleep, internal SRAM for DMA buffers (PSRAM cannot be a DMA
source), flash for constants.}
\label{tab:memmap}
\end{table}

\section{Controlling memory: bounded allocation}
\label{sec:memory}\index{Storage\_Size}

Embedded firmware usually wants allocation to be \emph{bounded} or absent. Ada
makes the size explicit rather than leaving it to a global heap:

\begin{itemize}
  \item \texttt{for T'Storage\_Size use 0;} on an access type forbids
        \texttt{new} of \texttt{T} -- a compile-time guarantee that pointer type
        never heap-allocates.
  \item A task can cap its own stack: \texttt{task body Worker} with
        \texttt{pragma Storage\_Size} (or the \texttt{Storage\_Size} aspect) sizes
        the primary stack precisely.
  \item The binder caps the secondary stacks (\texttt{-D<n>}, \texttt{-Q<n>}), and
        the \texttt{light-tasking} profile removes the heap and secondary stack
        entirely (a bump allocator only).
\end{itemize}

\begin{lstlisting}
   type Node_Ref is access Node with Storage_Size => 0;  -- new Node => compile error
   task Sampler with Storage_Size => 4 * 1024;           -- a 4 KB primary stack
\end{lstlisting}

Sizing memory in the source -- not discovering it when the heap is exhausted at
3 a.m.\ in the field -- is the embedded mindset the language supports directly.

\section{Interfacing with C, ROM, and assembly}
\label{sec:interfacing}\index{interfacing!with C}

A bare-metal stack still meets non-Ada code: a ROM routine, a C glue function, a
few instructions only assembly can express. Ada's interfacing pragmas make those
boundaries explicit and typed.

\textbf{Calling out} -- \texttt{pragma Import} (or the \texttt{Import} aspect)
with \texttt{Convention => C} binds an Ada declaration to an external symbol; no
header, no wrapper:

\begin{lstlisting}
   procedure Rom_Printf (Fmt : System.Address)
     with Import, Convention => C, External_Name => "esp_rom_printf";
\end{lstlisting}

\textbf{Being called} -- \texttt{pragma Export} publishes an Ada subprogram or
object under a C name so the boot assembly or glue can reach it (the runtime
exports its interrupt-dispatch entries and the per-core flags this way).
\texttt{Interfaces.C} (and \texttt{.Strings}, \texttt{.Pointers}) provides the
C-compatible scalar and array types for the data that crosses.

\textbf{Inline assembly} -- for the handful of operations with no Ada spelling
(a special register read, a cache or barrier instruction),
\texttt{System.Machine\_Code.Asm} embeds Xtensa instructions with typed operands;
the GDMA driver uses it for a memory barrier around descriptor hand-off. It is a
scalpel, not a habit: everything that can be Ada, is.

\section{The thread through it all}

One idea runs through this chapter and the last: in Ada the \emph{type} -- with
its layout, its range, its volatility, its contracts -- carries the intent, so the
compiler (or a cheap, switch-off-able run-time check) catches the mistake. Describe
a register, a value, or an interface correctly once, and every use of it afterwards
is safe by construction.
